\inicializarSubseccion{Descripción}
\tab Esta sección muestra cómo simulamos en MATLAB el campo eléctrico generado por un dipolo eléctrico, que no es más que dos cargas de igual magnitud pero signo opuesto separadas por una distancia $d$. La idea principal es calcular el campo de una serie de puntos del espacio y representarlo con flechas que indican hacia dónde apunta en cada lugar. El código permite modificar tanto la magnitud de las cargas como la separación entre ellas, lo que hace posible observar cómo cambia el campo ante distintas configuraciones. Esta simulación tiene una conexión directa con el reto ya que la dielectroforesis funciona precisamente exponiendo células a campos eléctricos no uniformes como los que genera este tipo de sistema.